\newpage
\section*{第22章：类型重构}
\addcontentsline{toc}{section}{第22章：类型重构}

\begin{taplsolutionentrybox}
\phantomsection
\label{sol:translator:22.2.3}
\textbf{22.2.3 解答}

记题中项为 \(t\)。要让 \((x\,z)(y\,z)\) 可赋型，\(x\,z\) 必须得到一个
函数，而 \(y\,z\) 必须得到该函数所需的参数。下面三组彼此不同的替换都满足
这一要求：
\[
\begin{array}{l}
\sigma_1=[X\mapsto\tapltype{Nat}\to\tapltype{Bool}\to\tapltype{Nat},
Y\mapsto\tapltype{Nat}\to\tapltype{Bool},Z\mapsto\tapltype{Nat}],\\
\sigma_2=[X\mapsto\tapltype{Bool}\to\tapltype{Nat}\to\tapltype{Bool},
Y\mapsto\tapltype{Bool}\to\tapltype{Nat},Z\mapsto\tapltype{Bool}],\\
\sigma_3=[X\mapsto Z\to A\to B,Y\mapsto Z\to A].
\end{array}
\]
相应结果类型分别是主类型的自然数实例、布尔实例，以及最一般的
\[
(Z\to A\to B)\to(Z\to A)\to Z\to B.
\]
第三组保留了 \(Z,A,B\)，因而比前两组更一般：再令
\((Z,A,B)=(\tapltype{Nat},\tapltype{Bool},\tapltype{Nat})\)，就得到第一组；
再令 \((Z,A,B)=(\tapltype{Bool},\tapltype{Nat},\tapltype{Bool})\)，就得到第二组。

\taplsolutionbacklink{ex:22.2.3}{返回练习22.2.3}
\end{taplsolutionentrybox}

\begin{taplsolutionentrybox}
\phantomsection
\label{sol:translator:22.3.3}
\textbf{22.3.3 解答}

先在上下文 \(x:X,y:Y,z:Z\) 中处理最内层的三个应用：
\begin{enumerate}
  \item \(x\,z\) 引入新鲜变量 \(A\)，生成约束 \(X=Z\to A\)，结果类型为 \(A\)；
  \item \(y\,z\) 引入新鲜变量 \(B\)，生成约束 \(Y=Z\to B\)，结果类型为 \(B\)；
  \item \((x\,z)(y\,z)\) 引入新鲜变量 \(C\)，生成约束 \(A=B\to C\)，结果类型为 \(C\)。
\end{enumerate}
因此，函数体 \((x\,z)(y\,z)\) 的暂定类型为 \(C\)。再由内向外应用
三次 CT-Abs，依次得到
\[
Z\to C,
\qquad Y\to Z\to C,
\qquad X\to Y\to Z\to C.
\]
于是整个项的暂定结果类型以及生成的新鲜变量集合和约束集可以取为
\[
\begin{aligned}
S&=X\to Y\to Z\to C,\\
\mathcal X&=\{A,B,C\},\\
C_0&=\{X=Z\to A,\ Y=Z\to B,\ A=B\to C\}.
\end{aligned}
\]
推导树中 \(x,z,y,z\) 四个变量出现分别由四次 CT-Var 得到；随后应用三次
CT-App 构造三个应用，再应用三次 CT-Abs 包回三层 lambda 抽象，最终得到
结论 \(\vdash t:S\mid_{\mathcal X}C_0\)。

\taplsolutionbacklink{ex:22.3.3}{返回练习22.3.3}
\end{taplsolutionentrybox}

\begin{taplsolutionentrybox}
\phantomsection
\label{sol:translator:22.5.2}
\textbf{22.5.2 解答}

合一\hyperref[sol:translator:22.3.3]{解答22.3.3}得到的约束：先由 \(A=B\to C\)，再代入
\(X=Z\to A\)、\(Y=Z\to B\)。可取主合一子
\[
[X\mapsto Z\to B\to C,
  Y\mapsto Z\to B,
  A\mapsto B\to C].
\]
把它作用于结果类型 \(S\)，得到主类型
\[
(Z\to B\to C)\to(Z\to B)\to Z\to C.
\]
变量名本身无关紧要；把 \(Z,B,C\) 一致地改名，仍是同一个主类型方案。

\taplsolutionbacklink{ex:22.5.2}{返回练习22.5.2}
\end{taplsolutionentrybox}

\begin{taplsolutionentrybox}
\phantomsection
\label{sol:translator:22.5.5}
\textbf{22.5.5 解答}

实现分为三层。第一层按\cref{fig:22-1}遍历项，返回暂定结果类型和约束；第二层
调用\cref{fig:22-2}的合一算法；第三层把主合一子作用于暂定结果类型。关键
不变量是：每次处理应用时，为其结果创建的类型变量都不曾出现过。无标注抽象要到
\taplsecref{22.6}{sec:22.6}才加入语言，本题仍沿用显式标注的 lambda 抽象。

书内已经依次给出实际 Rust 实现：\hyperref[fig:22-1]{图22-1后的类型与项定义}、
\hyperref[sol:author:22.3.10]{解答22.3.10中的约束生成器}和
\hyperref[sol:author:22.4.6]{解答22.4.6中的合一器}。把约束生成与合一串联起来、
再把所得替换作用于暂定结果类型的主类型入口如下：
\taplrustsupport{ch22-principal-type}

这三层合在一起，即构成题目要求的完整实现；这里不再重复刊载前两层源码。

同一份源码中的端到端测试覆盖题目给出的三组典型交互、自应用的出现检查失败、
多步替换的传递应用、显式标注与自动生成变量同名时的回归情形，以及布尔值、
自然数和条件项的约束生成。

\taplsolutionbacklink{ex:22.5.5}{返回练习22.5.5}
\end{taplsolutionentrybox}

\begin{taplsolutionentrybox}
\phantomsection
\label{sol:translator:22.5.7}
\textbf{22.5.7 解答}

本题只修改\cref{ex:22.5.5}的
\texttt{reconbase} 检查器，因此处理的仍是带显式形参类型标注的简单语言；
无标注抽象和 let 多态分别到\taplsecref{22.6}{sec:22.6}和
\taplsecref{22.7}{sec:22.7}再加入。

增量版本不再先收集整棵语法树的约束、最后统一求解，而是在处理每个语法构造时
立即合一新产生的约束。对外的入口仍返回整个项的主类型；递归辅助函数
\(\mathit{infer}\) 则返回一对 \((\sigma,T)\)：\(\sigma\) 是处理当前子项时
求得的替换，\(T\) 是在当前已知约束下得到的最一般类型；这里约定，返回的
\(T\) 已经应用了同一次调用求出的 \(\sigma\)。处理后续子项前，必须先把
\(\sigma\) 作用于上下文，使后续步骤能够使用刚刚获得的类型信息。

以应用 \(t_1\,t_2\) 为例，处理顺序可以明确写成：
\[
\begin{aligned}
  \mathit{infer}(\Gamma,t_1)&=(\sigma_1,T_1),\\
  \mathit{infer}(\sigma_1\Gamma,t_2)&=(\sigma_2,T_2).
\end{aligned}
\]
接着创建全新的结果类型变量 \(X\)，并求下列等式的最一般合一子 \(\sigma_3\)：
\[
  \sigma_2(T_1)=T_2\to X.
\]
左侧先应用 \(\sigma_2\)，是因为重构 \(t_2\) 时得到的新信息也可能约束
\(t_1\) 的类型。整个应用最终返回
\[
  (\sigma_3\circ\sigma_2\circ\sigma_1,\;\sigma_3(X)).
\]
复合顺序记录了信息产生的先后次序：后求得的替换必须作用于此前替换的结果。

这里的中间类型只是当前已处理约束下的最一般类型，后续约束仍可能使它更具体。
最一般性得以保留，是因为每次合一都选取最一般合一子：任何满足新约束的其他
替换，都可以写成这个最一般合一子与另一替换的复合。因此，把它复合到已有替换
时，只会施加满足新约束所必需的限制；处理到根节点时，所有约束都已纳入，所得
类型便是整个项的主类型。例如，
\((\lambda x:Y.x)\ 0\) 先为函数得到 \(Y\to Y\)，再由应用产生
\(Y\to Y=\tapltype{Nat}\to X\)；最一般合一子只令
\(Y=X=\tapltype{Nat}\)，于是整个应用的类型为 \(\tapltype{Nat}\)，无需重新
遍历函数体。

下面给出 \texttt{reconbase} 全部语法分支的增量实现。递归辅助函数负责逐层
处理项，并把替换返回给上一层，用来更新后续子项所处的上下文。最外层没有后续
子项，\texttt{principal\_type\_incremental} 因而只返回已经更新完毕的类型。
\taplrustsupport[breakable]{ch22-incremental-reconbase}

\taplsolutionbacklink{ex:22.5.7}{返回练习22.5.7}
\end{taplsolutionentrybox}

\begin{taplsolutionentrybox}
\phantomsection
\label{sol:translator:22.7.1}
\textbf{22.7.1 解答}

在\cref{ex:22.5.7}的增量重构器上增加类型方案
\(\forall X_1\ldots X_n.T\)。处理 let 时，先重构右侧并把当前替换作用于
上下文；计算右侧类型中自由、但上下文中不自由的变量，把它们泛化后绑定给
let 变量。处理变量时，每次查询类型方案都以全新的类型变量实例化所有量化变量。

若语言包含引用等副作用，泛化步骤还必须实行值限制：只有右侧是句法值时才泛化；
其他右侧以不量化的单态方案加入上下文。下面先给出在本节中新增加的类型方案、
泛化、实例化和 let 分支等完整实现，再给出计算主类型的入口。
\taplrustsupport[breakable]{ch22-algorithm-w-complete}

\taplsolutionbacklink{ex:22.7.1}{返回练习22.7.1}
\end{taplsolutionentrybox}
